% 319_Proton_Torus_De_ch.tex
% Das Proton als schwingender Torus:
% Geometrische Grundlagen der Massenemergenz in der FFGFT
% Dok. 319, August 2026

\section*{Zusammenfassung}

Das Proton wird in der FFGFT nicht als Ansammlung von Quarks und
Gluonen in einem Volumen beschrieben, sondern als \emph{schwingender
Torus} $T^4/\mathbb{Z}_3$. Quarks und Gluonen sind verschiedene
Modentypen dieser Struktur. Die Protonmasse ist die Schwingungsenergie
des Torus. Dieses Dokument entwickelt das geometrische Bild, verankert
es im Korpus und markiert alle Aussagen nach dem epistemischen Standard
[K]/[S]/[SETZUNG].

\begin{tcolorbox}[colback=blue!5!white,colframe=blue!60!black,
  title=Zur Lesart: Innen-Perspektive]
\textbf{[K]} Wir beschreiben den Proton-Torus \emph{von innen} --
als Teil desselben $T^4$-Systems (Dok.~248). Es gibt keine
Außenperspektive. Was ``Einschluss'' oder ``Freiheit'' bedeutet,
ist relational: es beschreibt die Relation zwischen dem Beobachter
und dem jeweiligen topologischen System, nicht eine intrinsische
Eigenschaft des beobachteten Felds.
\end{tcolorbox}

\section{Das geometrische Grundbild}

\textbf{[K]} Die FFGFT geht von einem kompakten 4D-Torus $T^4$ mit
$\mathbb{Z}_3$-Orbifold-Symmetrie aus. Die fundamentale Relation ist
(Dok.~003):
\begin{equation}
  \tilde{T} \cdot m = 1
\end{equation}
Intrinsische Zeit und Masse sind dual. Das hat unmittelbare Konsequenzen
für alle Moden auf dem Torus.

\textbf{[K]} Die Trialität des D4-Gitters ist in der Gittergeometrie
eingebaut, nicht gesetzt (Dok.~314): Die Ordnung-3-Elemente in
$\text{Aut}(D4)$ sind \emph{halbzahlig} -- die Klebevektor-Mischung
$(\tfrac{1}{2},\tfrac{1}{2},\tfrac{1}{2},\tfrac{1}{2})$ geht ein. Die
$T^4/\mathbb{Z}_3$-Struktur mit genau 9 Fixpunkten folgt aus der
D4-Geometrie. Diese halbzahligen Elemente sind der geometrische Ursprung
der fermionischen Statistik der Quark-Moden.

\section{Zwei Modentypen: Quarks und Gluonen}

\textbf{[K]} Die kompakte $\mathbb{Z}_3$-Faser trägt das diskrete,
massentragende Spektrum -- nicht die ausgerollte Achse (Dok.~285):
\begin{quote}
``Das diskrete, massentragende Spektrum sitzt auf der kompakten
$\mathbb{Z}_3$-Faser, nicht auf der ausgerollten Achse.''
\end{quote}

\begin{table}[h]
\centering
\caption{Quark- und Gluon-Moden auf dem Torus [K]}
\label{tab:moden}
\begin{tabular}{lllll}
\toprule
Mode & Spektrum & Masse & Eigenzeit $\tilde{T}$ & Dok. \\
\midrule
Quark-Moden & diskret, $\mathbb{Z}_3$-Faser
  & $m > 0$ & $\tilde{T} < \infty$ & 006, 285, 314 \\
Gluon-Moden & kontinuierlich, lokal flach
  & $m = 0$ & $\tilde{T} = \infty$ & 003, 270 \\
\bottomrule
\end{tabular}
\end{table}

\textbf{[K]} Die Gluonen sind masselos ($m = 0$). Damit gilt
$\tilde{T} = \infty$ -- sie vollziehen keinen Messvorgang (Dok.~270).
Lokal gilt Flachheit: „im räumlichen Limes die lokale Flachheit"
(Dok.~270). Das bedeutet: lokal sieht das masselose Gluon $\mathbb{R}^3$,
aber das ist ein lokaler Homöomorphismus, keine globale Aussage.

\section{Die relationale Natur von Einschluss und Freiheit}

\textbf{[K]} Das Ausrollen ist ein Messvorgang des eingebetteten
Beobachters (Dok.~270): „Der Messvorgang wählt einen Kreis als Zeit
und rollt ihn auf; ein physikalischer Akt, kein mathematischer
Grenzübergang."

\textbf{[K]} Wir (Beobachter) sind Teil des $T^4$-Systems -- nicht
außerhalb (Dok.~248): „Wir sind Teil des Systems. Wir beschreiben
immer von innen, nie von außen."

Daraus folgt eine fundamentale Einsicht für das Photon-Gluon-Verhältnis:

\begin{keyresult}[Einschluss ist relational {[K/S]}]
Photon und Gluon sind beide masselos ($m=0$), beide zeitlos
($\tilde{T}=\infty$), beide sehen lokal $\mathbb{R}^3$.
\textbf{Beide tragen Randbedingungen.}

Der Unterschied ist nicht intrinsisch, sondern \emph{relational}:
\begin{itemize}
  \item \textbf{Photon:} Wir (Beobachter) sind Teil desselben
    $T^4$-Systems. Das Photon bewohnt denselben topologischen Sektor.
    Seine Randbedingung ist unsere Randbedingung. Es erscheint uns
    daher als frei propagierend.
  \item \textbf{Gluon:} Wir stehen der $\mathbb{Z}_3$-Unterstruktur
    des Proton-Torus von außen gegenüber. Das Gluon bewohnt einen
    anderen topologischen Sektor. Es erscheint uns als gebunden.
\end{itemize}
``Einschluss'' ist keine Eigenschaft des Gluons allein -- es ist
eine Eigenschaft der Relation Beobachter $\leftrightarrow$ System
(Dok.~248, [K]). \textbf{[S]} Die vollständige topologische Herleitung
warum Photon und Gluon verschiedene Sektoren bewohnen, ist noch offen.
\end{keyresult}

\section{Masselose Feldmoden: lokal pure Energie}

\textbf{[K]} Für masselose Felder ($m=0$) ist $\tilde{T}\cdot m=1$
nicht definiert -- nicht als Grenzwert, sondern exakt: Auf
Nullgeodäten gilt $d\tau = 0$; entlang eines Lichtwegs vergeht
keine Eigenzeit. Es gibt kein Ruhesystem (Dok.~312):
\begin{quote}
``Masselos bedeutet $v=c$ exakt -- die Sättigung der Ungleichung
$v\le c$. Es gibt kein $c^*(\lambda)$; $dc/d\lambda\equiv 0$.''
\end{quote}

\textbf{[K]} Was lokal bleibt, ist ausschließlich $E = pc$ -- reine
Energie (Dok.~134, 312). Das gilt für Photon-Feldmoden und
Gluon-Feldmoden \emph{gleichermaßen}. Aus lokaler Sicht ist jede
masselose Feldmode pure Energie. Es gibt kein ``Gluon hier'' oder
``Photon dort'' in einem absoluten Sinn -- nur Energiedichte im Raum.

\textbf{[K]} Auch ``Teilchen'' sind in der FFGFT nicht fundamental
(Dok.~049, 156): ``Es gibt überhaupt keine Teilchen! Was wir Teilchen
nennen, sind verschiedene Anregungsmuster im einzigen Feld
$\delta m(x,t)$.'' Photon und Gluon sind Feldmoden -- keine Objekte.

\textbf{[S]} Der Unterschied Photon/Gluon liegt nicht in der lokalen
Beschreibung (beide: $E=pc$, beide: $d\tau=0$, beide: lokal flach),
sondern in der \emph{globalen topologischen Einbettung} der Feldmode:

\begin{center}
\begin{tabular}{lll}
\toprule
Feldmode & Topologie & Globales Verhalten \\
\midrule
Photon & $U(1)$-Fluss, ausgedehnt & propagiert in $T^4$ (unser System) \\
Gluon & $\mathbb{Z}_3/SU(3)$-Verschlingung & stehende Welle im Proton-Torus \\
\bottomrule
\end{tabular}
\end{center}

\textbf{[S]} Das Gluonfeld ist nicht lokalisierbar: Als
nicht-Abelsche $SU(3)$-Feldmode (Dok.~145: Farb-Antifarb-Verschlingung
dreier Flussfäden) wechselwirkt sie mit sich selbst. Es gibt kein
abgrenzbares ``Gluon-Objekt'' -- nur eine globale Feldkonfiguration
des Proton-Torus. Die acht Gluonmoden entsprechen den acht
nicht-trivialen Verschlingungszuständen der drei Flussfäden =
acht Generatoren von $SU(3)_c$ (Dok.~145, [S]).

\section{Die topologische Randbedingung als Confinement-Mechanismus}

\textbf{[K]} Auf dem geschlossenen Torus muss eine Welle eine
ganzzahlige Phasenverschiebung über den Umfang aufweisen:
\begin{equation}
  \oint k \, \mathrm{d}l = 2\pi n, \quad n \in \mathbb{Z}
\end{equation}

\textbf{[K]} Die $\mathbb{Z}_3$-Wirkung auf $L^2(T^4)$ permutiert
die Moden-Orbits zyklisch; die Permutationsmatrix hat Eigenwerte
$\{1, \omega, \omega^2\}$, Phasen $0^\circ$, $+120^\circ$,
$-120^\circ$ -- strukturell exakt die Zirkulant-Diagonalisierung
der $\mathbb{C}^3$-Faser (Dok.~314). Das erzwingt diskrete Moden
ohne Wände.

\textbf{[S]} Der Schritt von dieser Modenstruktur zur vollen
QCD-Confinement-Dynamik ist noch nicht formal ausgeführt.

\section{Die Dualität $\tilde{T} \cdot m = 1$ im Proton}

\textbf{[K]} Die Vollständige Energierelation
$E^2 = (m_0 c^2)^2 + (pc)^2$ (Dok.~318) zeigt: masselose Gluonen
tragen trotz $m_0 = 0$ die Energie $E = pc$ -- sie tragen $\sim 99\,\%$
der Protonmasse als Feldenergie.

\textbf{[K]} Das ist keine Paradoxie. Eine stehende Welle hat keine
Eigenzeit ($\tilde{T} = \infty$), aber ihre Energie $E = \hbar\omega$
ist real und messbar. Die Dualität $\tilde{T} \cdot m = 1$ besagt:
Zeitlosigkeit und Masselosigkeit sind äquivalent. Die
\emph{Systemenergie} des gebundenen Zustands -- nicht die Eigenenergie
der einzelnen Mode -- ist die Protonmasse.

\begin{table}[h]
\centering
\caption{Zeitbehaftete vs. zeitlose Schwingung [K]}
\begin{tabular}{lll}
\toprule
Eigenschaft & Quark (zeitbehaftet) & Gluon (zeitlos) \\
\midrule
Eigenzeit & $\tilde{T} > 0$ & $\tilde{T} = \infty$ \\
Energie & Ruhemasse $+$ Bewegung & Reine Feldenergie $E = pc$ \\
Spektrum & Diskret ($\mathbb{Z}_3$-Faser) & Lokal kontinuierlich \\
Statistik & Fermion (D4-Trialität [K]) & Boson \\
\bottomrule
\end{tabular}
\end{table}

\section{Spin-Statistik aus D4-Trialität}

\textbf{[K]} Die fermionische Statistik der Quark-Moden folgt aus
der D4-Gitter-Geometrie (Dok.~314): Die 80 Ordnung-3-Elemente in
$\text{Aut}(D4)$ zerfallen in drei Klassen; die Klasse mit Spur $-2$
ist \emph{halbzahlig} und fixpunktfrei in $\mathbb{R}^4$. Diese
Halbzahligkeit ist die geometrische Grundlage der Fermion-Statistik.

\textbf{[S]} Der Entwurf-Ansatz „$(n_\theta + n_\phi)$ ungerade
$\to$ Fermion" trifft das Ergebnis, aber nicht den Weg. Korrekt ist:
Die D4-Trialität liefert halbzahlige Darstellungen, unabhängig von der
Parität der Wicklungszahlen.

\section{Farbladung und $\mathbb{Z}_3$-Struktur}

\textbf{[K]} Die $\mathbb{Z}_3$-Trialität des D4-Gitters ist
eingebaut, nicht gesetzt (Dok.~314). $\mathbb{Z}_3$ ist das Zentrum
von $SU(3)$.

\textbf{[S]} Ob und wie daraus die volle $SU(3)_c$-Eichstruktur
emergiert, ist im Korpus noch nicht formal hergeleitet.

\textbf{[S]} Dok.~049 formuliert: „Nodes cannot exist alone, must
form neutral combinations." Das ist eine Formulierung des Confinements,
aber kein Beweis der $SU(3)_c$-Emergenz.

\begin{tcolorbox}[colback=orange!8!white,colframe=orange!60!black,
  title=Status der Farbladungs-These {[S]}]
Die Aussage ``Farbladung ist effektiv, nicht fundamental'' ist ein
geometrisches Bild \textbf{[S]}, keine bewiesene Ableitung. Der
$\mathbb{Z}_3$-Zusammenhang [K] ist ein struktureller Hinweis auf
die $SU(3)_c$-Struktur -- die vollständige geometrische Herleitung
bleibt offen.
\end{tcolorbox}

\section{$\Lambda_{\text{QCD}}$ und minimale Frequenz}

\textbf{[S]} Die minimale Frequenz auf einem Torus mit Radius
$R \approx 1$~fm:
\begin{equation}
  \omega_{\min} = \frac{1}{R} \approx 197~\text{MeV}
  \approx \Lambda_{\text{QCD}}
\end{equation}
($1~\text{fm}^{-1} = 197{,}3$~MeV in natürlichen Einheiten).
Konsistent mit Dok.~041 ($\Lambda_{\text{QCD}} = v\cdot\xi^{1/3}
\approx 200$~MeV), aber keine vollständige Herleitung.

\section{Das Proton als Ganzes}

\textbf{[K]} Das Proton ist kein Behälter mit Teilchen, sondern eine
\emph{Systemgröße}: Die Quark-Moden (diskret, $\mathbb{Z}_3$-Faser,
massebehaftet) definieren die Randbedingung. Die Gluon-Moden (lokal
kontinuierlich, masselos) füllen das System mit Feldenergie. Die
Protonmasse ist die Gesamtenergie dieses gekoppelten Systems.

\textbf{[S]} Das geometrische Bild „Quarks am Rand, Gluonen im
Volumen" ist konzeptuell belastbar, aber noch keine formale Ableitung
aus der Torus-Geometrie.

\section{Bezug zum Korpus}

\begin{itemize}
  \item \textbf{Dok.~003}: $\tilde{T}\cdot m=1$, $D_\text{frak}=2{,}94$,
    $K_\text{frak}$ [K]
  \item \textbf{Dok.~248}: Eingebetteter Beobachter, Innen-Perspektive [K]
  \item \textbf{Dok.~270}: Messvorgang als Ausrollen, lokale Flachheit [K]
  \item \textbf{Dok.~285}: Diskretes Spektrum auf $\mathbb{Z}_3$-Faser [K]
  \item \textbf{Dok.~311}: Eingerollt/ausgerollt, drei Wege [K]
  \item \textbf{Dok.~314}: D4-Trialität, halbzahlig, $T^4/\mathbb{Z}_3$
    aus D4 [K]
  \item \textbf{Dok.~318}: Proton als offene Brücke, $E^2=(m_0c^2)^2+(pc)^2$,
    Geltungsbereich [K]
  \item \textbf{Dok.~041}: $\Lambda_\text{QCD}\approx 200$~MeV [S]
  \item \textbf{Dok.~049}: Confinement-Formulierung [S]
\end{itemize}

\section{Offene Fragen [S]}

\begin{enumerate}
  \item \textbf{$SU(3)_c$-Emergenz aus $\mathbb{Z}_3$-Trialität.}
    Der strukturelle Hinweis [K] ist belegt. Die vollständige
    Eichstruktur-Herleitung fehlt.
  \item \textbf{Warum Photon und Gluon verschiedene Sektoren.}
    Die relationale Erklärung (Dok.~248) ist [K]. Die formale
    topologische Herleitung bleibt [S].
  \item \textbf{$\Lambda_\text{QCD}$ aus Torus-Geometrie.}
    Dok.~041: Ordnungsgröße [S]. Vollständige Herleitung offen.
  \item \textbf{RGE-Brücke $\alpha_s(m_\tau)\to\alpha_s(M_Z)$.}
    Dok.~160: $\alpha_s(m_\tau)=3\xi^{1/4}$ [K]. Brücke offen [S].
  \item \textbf{Formale Randbedingungen auf $T^4/\mathbb{Z}_3$.}
    Exakte Spektraltheorie noch auszuarbeiten [S].
\end{enumerate}
